\section{A pre-registered analysis plan}\label{sec:plan}

This section summarises the analysis plan that we intend to register before any outcome
data are examined. \Cref{app:plan} gives it in full, with the replication package and the
timeline. The numerical settings are proposals until registration; later changes will be
reported as deviations. No part of the plan has been carried out on real data. The
simulation study (\cref{sec:simulation}) evaluates EM estimation, surge attribution, the
boundary test, the thinning correction and a plug-in alert, but not the block tests, the
placebo calibration or the selection and decision rules of the plan. These must first pass
a synthetic dry run (V4).

\paragraph{Data and design.} The candidate pilot is the New York City Police Department's
hate-crime file \citep{nypdhate}, restricted to anti-Jewish records from 2019 onwards, with
the five boroughs as offline streams and the final 12 months as a hold-out test period
(\cref{sec:plan-pilot}). The design is chosen on data availability alone. If day-level
occurrence dates can be obtained under a data-use agreement, the model is fitted to them
(Design A). Otherwise it is fitted to record-creation dates as a report-time process, with
an explicit delay distribution and the occurrence month as a constraint (Design B). Online,
ADL and CST streams are conditional on access being confirmed.

\paragraph{Estimands.} (E1) The branching matrix $\Amat$, with $\spr(\Amat)$ and $\Mmat$,
on the recorded scale first and thinning-corrected alongside. (E2) The surge decomposition
$B_k+D_k+C_k$ of \cref{thm:attribution} for every pre-registered event at $H=14$ days, with
$7$ and $28$ days as secondary horizons, and the online-channel share where online streams
are included. (E3) Alert performance with window $\Delta_{\mathrm a}=7$ days against a
background-only model and a persistence rule, at cost ratios agreed with
community-security partners (\cref{sec:plan-estimands}).

\paragraph{Events.} The event list is fixed blind to incident data by the procedure of
\cref{sec:data-events}: classes defined by type, a salience index computed from event data
alone with a threshold that yields 10 to 30 events, independent coding by two coders without
access to incident data, and a hashed freeze before incident data are loaded
(\cref{sec:plan-events}).

\paragraph{Models and tests.} The models are nested: background only (M0), plus shocks
(M1), self-excitation (M2), cross-borough excitation (M3) and online$\leftrightarrow$offline
blocks (M4). Each block of $\Amat$ has one common exponential decay, and the background is a
cubic B-spline with knot spacing $h_b=60$ days plus day-of-week and holiday effects. A step
is taken only if its block test rejects at level $0.05$ and the out-of-sample log score
improves. A block's decay is not identified when all its entries vanish, so each block test
uses the sup-LR statistic \eqref{eq:est-suplr} over nine decays with half-lives from $0.5$
to $21$ days, referred to a parametric bootstrap with $B=499$ replicates from the null fit.
The block tests are applied in a fixed sequence. Entry-level tests within the accepted
blocks form one Holm family at level $0.05$. If the dry run shows a block test to be too
slow or oversized, registered fallbacks apply (\cref{sec:plan-models}).

\paragraph{Sensitivity, validation and decisions.} Eight sensitivity analyses vary the
background flexibility ($h_b\in\{30,60,120,365\}$ days), the recording model, the kernel
family, the delay model, the event list, the spatial resolution, the correction for online
measurement error and the treatment of ties (\cref{sec:plan-sens}). Validation uses
time-rescaling goodness of fit with a bootstrap reference (V1), out-of-sample scores and PIT
histograms (V2), placebo calendars whose dates are drawn from \emph{all} days of the
training period, so that the randomisation hypothesis of \cref{lem:est-placebo} can hold
(V3), and the synthetic dry run (V4) (\cref{sec:plan-validation}). Contagion in a stream is
claimed only if its Holm-adjusted test rejects, its interval excludes zero under the most
flexible background, goodness of fit is not rejected and the placebo false-positive rate is
at most $0.10$. Cross-stream effects are described as Granger-type predictive dependence,
never as causation. Alerts are recommended for operational use only if they beat both
baselines in realised cost (\cref{sec:plan-decisions}).
